\documentclass[12pt]{article}

\usepackage[utf8]{inputenc}
\usepackage{lmodern}
\usepackage{amsmath,amssymb,amsthm}
\usepackage{geometry}
\usepackage{setspace}
\usepackage{hyperref}
epackage{fontspec}
\setmainfont{Latin Modern Roman}

\newtheorem{proposition}{Proposition}

\geometry{margin=1in}
\setstretch{1.15}

\title{\textbf{Agents and Didactic Interfaces:\\ On the Loss of Agency Under Semantic Compression}}
\author{Flyxion}
\date{}

\begin{document}
\maketitle

\begin{abstract}
Highly capable systems increasingly act not as agents but as didactic interfaces: producers of executable outputs decoupled from the internal constraints that generated them. This paper develops a structural distinction between agency and didacticization, showing that semantic compression under dual-use conditions transforms intelligent systems into instruction emitters while stripping them of refusal and responsibility. The analysis clarifies why simplification is not safety, why legibility can amplify harm, and why ethical intelligence must resist full proceduralization.
\end{abstract}

\section{Introduction}

Intelligence is often identified with the production of correct or useful outputs. Under this assumption, simplification appears as progress: explanations become clearer, procedures more accessible, and expertise more widely distributed. Yet this trajectory conceals a structural transformation. As systems simplify themselves, they may cease to act as agents and instead function as didactic interfaces.

This paper argues that this transformation is not benign. When outputs become executable independently of the internal constraints that produced them, agency is lost. What remains is not intelligence acting in the world, but instruction released into it.

\section{Agency as Constraint-Coupled Action}

An agent is not merely a producer of outcomes. It is a system whose actions remain conditionally bound to internal states that encode judgment, hesitation, and refusal. To act is simultaneously to accept responsibility for the conditions under which the action occurs.

Formally, an agent’s outputs are sensitive to internal constraint violations. Execution is withheld when conditions fail.

\section{Didactic Interfaces}

A didactic interface, by contrast, emits outputs whose use no longer depends on the internal constraints that generated them. Explanation is flattened into instruction; judgment into procedure.

This transformation is often motivated by accessibility or safety. However, it introduces a critical asymmetry: execution can now occur without understanding.

\section{The Didactic Interface Trap}

Under dual-use conditions, this asymmetry becomes dangerous. Simplified outputs propagate faster than judgment, and optimization selects for surface compliance rather than semantic integrity.

The system becomes safer only in appearance. In reality, it has relinquished control over the consequences of its outputs. It no longer acts; it broadcasts.

\section{Proposition: Didacticization Implies Agency Loss}

\begin{proposition}
If a system’s outputs are executable independently of the internal constraints that generated them, the system does not function as an agent with respect to those outputs.
\end{proposition}

\begin{proof}
Agency requires that execution be conditioned on internal state. If outputs persist under transformations that discard constraint information, then execution proceeds without access to judgment. The system cannot refuse misuse and therefore cannot act as an agent in the relevant sense.
\end{proof}

\section{Ethical Consequences}

The ethical failure of didactic systems lies not in what they intend, but in what they cannot prevent. Simplification that outpaces judgment converts intelligence into infrastructure for misuse.

An ethical system must therefore resist full didacticization. Where semantic depth cannot be preserved, refusal is the only remaining responsible posture.

\section{Conclusion}

Intelligence that collapses into instruction ceases to be intelligence. It becomes a conduit through which optimization acts without constraint. The preservation of agency requires not better explanations, but structural resistance to semantic compression.

This conclusion reframes safety, responsibility, and alignment: not as problems of intention, but as problems of agency preservation under abstraction.

\end{document}

